\documentclass[11pt]{article}
\usepackage[margin=1in]{geometry}
\usepackage[T1]{fontenc}
\usepackage[utf8]{inputenc}
\usepackage{lmodern}
\usepackage{amsmath,amssymb,amsthm,mathtools}
\usepackage{booktabs,tabularx,longtable,array}
\usepackage{xltabular}% tabularx-style tables that can break across pages
\usepackage{enumitem}
\usepackage[dvipsnames]{xcolor}
\usepackage[noadjust]{cite}% sorts and compresses numeric citations; noadjust: no automatic space before \cite (it misfires after "(")
\usepackage[colorlinks=true,linkcolor=Blue,citecolor=Blue,urlcolor=Blue]{hyperref}

% ---------------------------------------------------------------------
% Shared notation
\newcommand{\R}{\mathbb{R}}
\newcommand{\E}{\mathbb{E}}
\renewcommand{\P}{\mathbb{P}}
\newcommand{\one}{\mathbf{1}}
\newcommand{\D}{\mathcal{D}}
\newcommand{\Lgen}{\mathcal{L}}
\DeclareMathOperator{\Var}{Var}
\DeclareMathOperator{\Cov}{Cov}
\DeclareMathOperator{\Ent}{Ent}
\DeclareMathOperator{\KL}{KL}
\DeclareMathOperator{\Tr}{Tr}
\DeclareMathOperator{\ran}{ran}
\newcommand{\norm}[1]{\lVert #1\rVert}
\newcommand{\ip}[2]{\langle #1,#2\rangle}
% Locators.  \ploc{compiled locator}{TeX file:line} cites the preprint (or a
% companion).  Only the compiled locator (section, equation, result, page of the
% PDF) is printed.  The TeX file:line stays in the source to help checkers find
% the passage and never appears in the PDF.  \sloc{key}{locator} cites a source.
\newcommand{\ploc}[2]{#1}
\newcommand{\sloc}[2]{\cite[#2]{#1}}
% Relationship types and flags, as used in the summary table and the Summary
% paragraphs: \rel{same technique}; \attrgap and \overstated tag a Cited?
% cell (precedence_stats.py checks them against the JSON).
\newcommand{\rel}[1]{\textit{#1}}
\newcommand{\attrgap}{\textbf{attribution gap}}
\newcommand{\overstated}{\textbf{overstated novelty}}
% Uniform list and table styling across all findings.
\setlist{itemsep=2pt,topsep=3pt,parsep=0pt,leftmargin=1.6em}
\setlist[enumerate]{label=\arabic*.}
\setlist[itemize]{label=\textbullet}
\renewcommand{\arraystretch}{1.15}
% Findings are numbered by LaTeX, so nobody types an ID.  In the main file,
% \findinggroup{F} after the Findings \section sets the ID prefix and restarts
% the count; IDs follow the order of the \input lines (by magnitude).  Each
% fragment opens with
%   \finding{<title>}{<slug>}   or   \finding[<short title>]{<title>}{<slug>}
% which prints an unnumbered subsection "F3. <title>", adds "F3. <title>" (or
% the short title) to the table of contents and PDF bookmarks, and does
% \label{f:<slug>}, so \ref{f:<slug>} prints the ID "F3" as a link.  The
% heading carries no LaTeX section number, so a finding never shows two
% numbers (no "3.2 F2").  Refer to another finding as finding~\ref{f:<slug>}.
% Before any \findinggroup the IDs print as ?1, ?2, ... (as in a standalone
% test of one fragment).
\makeatletter
\newcounter{finding}
\newcommand{\finding@prefix}{?}
\renewcommand{\thefinding}{\finding@prefix\arabic{finding}}
\newcommand{\findinggroup}[1]{\gdef\finding@prefix{#1}\setcounter{finding}{0}}
\newcommand{\finding}[3][]{%
  \stepcounter{finding}%
  \protected@edef\finding@id{\thefinding}% fix the ID now: "F3", not a counter reference
  \edef\finding@next{\noexpand\finding@sect{\finding@id}}%
  \finding@next{#1}{#2}{#3}}
% #1 ID (already expanded), #2 short title (may be empty), #3 title, #4 slug.
\newcommand{\finding@sect}[4]{%
  \subsection*{#1. #3}% hyperref anchors \subsection* itself, just above the heading
  \if\relax\detokenize{#2}\relax
    \addcontentsline{toc}{subsection}{#1. #3}%
  \else
    \addcontentsline{toc}{subsection}{#1. #2}%
  \fi
  \def\@currentlabel{#1}% \ref prints the ID; hyperref links it to this heading
  \label{f:#4}}
\makeatother
% ---------------------------------------------------------------------
\space%

\title{Missing citations in\\ ``The Unique Games Theorem''}
\date{October 16, 2026}

\begin{document}
\maketitle

\section*{Summary of findings}

Each row is a citation the preprint is missing: earlier work that already
contains part of the preprint, which the preprint does not cite where that
part appears. ``Cited?'' names the earlier work and says what, if anything, the
preprint cites it for. The exception is a row flagged \emph{overstated
novelty}: there the preprint cites the work, but its abstract or framing
presents the content as new. Terms are defined in Section~\ref{sec:about};
centrality is a best-effort judgment. Locations use the preprint's own
numbering and printed pages.

% Each row starts "\ref{f:<finding>} & \rel{<relationship>} & <centrality> &",
% in that order: \ref prints the finding's ID (F1, F2, ...) as a link, and
% precedence_stats.py reads those three cells. Put \attrgap and \overstated in
% the Cited? cell of a row exactly when the JSON sets the flag. Never type an
% ID: LaTeX assigns them from the main file's order (by magnitude).
\begin{center}\footnotesize
\begin{xltabular}{\linewidth}{@{}l >{\raggedright\arraybackslash}p{0.11\linewidth} >{\raggedright\arraybackslash}p{0.10\linewidth} >{\raggedright\arraybackslash}X >{\raggedright\arraybackslash}p{0.17\linewidth} >{\raggedright\arraybackslash}p{0.24\linewidth}@{}}
\toprule
ID & Relationship & Centrality & Technique or result & In the preprint & Cited?\\
\midrule
\endhead
\ref{f:erasure-covering-argument} & \rel{same technique} & central &
Erasure argument: re-randomize the function on the structured pieces already found, so that concentration, a union bound and the inverse theorem give a covering fraction &
Lemma~5.5 (pp.~26--27) &
Dinur, Khot, Kindler, Minzer, Safra 2016 \cite{DinurEtAl2016-Towards2to1-ECCC} (journal version \cite{DinurEtAl2018-Towards2to1}): cited for other content; \attrgap\\
\addlinespace
\ref{f:zoom-in-advice-zoom-out-extension} & \rel{adapted} & central &
Decoding with a shared zoom-in sent as advice, a zoom-out each prover finds in its own table, and a uniformly random extension as the answer &
Definition~5.2 (p.~24); Lemma~5.5 (p.~26); \S5.6 (p.~30) &
Dinur, Khot, Kindler, Minzer, Safra 2016 \cite{DinurEtAl2016-Towards2to1-ECCC} (journal version \cite{DinurEtAl2018-Towards2to1}): cited for other content; Khot, Minzer, Safra 2017 \cite{KhotMinzerSafra2017-IndependentSets2to2Grassmann}: credited in part; \attrgap\\
\addlinespace
\ref{f:folded-vertex-classes} & \rel{fragment} & central &
A folded constraint graph written out explicitly: question--table pairs identified across outer questions and intercept translates, one representative per class, an unfolding rule, aggregated weighted constraints &
\S\S4.3--4.4 (pp.~20--21) &
Dinur, Khot, Kindler, Minzer, Safra 2016 \cite{DinurEtAl2016-Towards2to1-ECCC} (journal version \cite{DinurEtAl2018-Towards2to1}): cited for other content; Khot, Safra 2013 \cite{KhotSafra2013-TwoProverStrongSoundness} and Bellare, Goldreich, Sudan 1998 \cite{BellareGoldreichSudan1998-FreeBitsPCPs}: credited for the two underlying principles\\
\addlinespace
\ref{f:kernel-clustering-khot-naor} & \rel{same result} & peripheral &
The identity-target kernel clustering problem and its rounding constant $(8\pi/9)(1-1/k)$ &
\S8.3, ``Kernel clustering'' (p.~43) &
Khot, Naor 2008 \cite{KhotNaor2008-ApproximateKernelClustering}: not cited (credited only through the same-release companion)\\
\addlinespace
\ref{f:correlation-clustering-via-multicut} & \rel{same route} & peripheral &
Correlation Clustering hardness drawn from Multicut hardness through their known equivalence &
Corollary~8.3, fourth conclusion (pp.~42--43) &
Chawla, Krauthgamer, Kumar, Rabani, Sivakumar 2006 \cite{ChawlaEtAl2006-HardnessMulticutSparsestCut}: cited for other content; Demaine, Emanuel, Fiat, Immorlica 2006 \cite{DemaineEtAl2006-CorrelationClusteringGeneralWeightedGraphs}: credited for the transformation\\
\addlinespace
\ref{f:earlier-unique-games-hardness-points} & \rel{weaker result} & peripheral &
Earlier NP-hardness of Unique Games at other parameter points, and conditional proofs under other hypotheses &
Theorem~1.1, \S\S1.2--1.3 (pp.~2--3) &
O'Donnell, Wright 2012 \cite{ODonnellWright2012-NewPointNPHardnessUniqueGames}: not cited; H{\aa}stad, Huang, Manokaran, O'Donnell, Wright 2017 \cite{HastadEtAl2015-ImprovedNPInapprox2LinearEquations}: not cited; Bhangale, Khot 2019 \cite{BhangaleKhot2019-UGHardnessToNPHardnessByLosingHalf}: not cited; Martinsson 2024 \cite{Martinsson2024-NPHardnessApproximationCurveMax2Lin2}: not cited; Raghavendra, Steurer 2010 \cite{RaghavendraSteurer2010-GraphExpansionUGC}: not cited; Eldan, Moshkovitz 2020 \cite{EldanMoshkovitz2020-NonUniqueToBooleanUniqueGames}: not cited; Wiman 2018 \cite{NL-Wiman2018}: not cited; Feige, Reichman 2004 \cite{NL-FeigeReichman2004}: not cited; H{\aa}stad 2001 \cite{Hastad2001-SomeOptimalInapproximability}: cited for other content\\
\addlinespace
\ref{f:strong-left-ordering-composition} & \rel{weaker result} & peripheral &
Unconditional Maximum Acyclic Subgraph hardness within $2/3+\epsilon$ from a strong-left Unique Games theorem, a weaker earlier version of the ordering threshold &
Lemma~A.1, Proposition~A.2, Corollary~8.2 (pp.~41--54); \S8.2 (p.~42) &
Bhangale, Khot 2019 \cite{BhangaleKhot2019-UGHardnessToNPHardnessByLosingHalf}: not cited\\
\bottomrule
\end{xltabular}
\end{center}

\clearpage
\space%

\tableofcontents

\section{About this document}\label{sec:about}

\paragraph{The preprint.}
The preprint is \emph{The Unique Games Theorem}, dated September 23, 2026, with author line
``OpenAI'' \cite{OAI-Preprint}. It belongs to the OpenAI Math Release of
September and October 2026, a collection of machine-generated mathematics
manuscripts published on GitHub. It proves the Unique Games Conjecture: for every
fixed $\varepsilon,\delta\in(0,1/2)$ it gives a deterministic polynomial-time
reduction from 3SAT to explicit unweighted bipartite Unique Games instances over
a fixed binary alphabet with translation constraints, with completeness at least
$1-\varepsilon$ and soundness at most $\delta$.
Section~\ref{sec:overview} summarizes the paper and fixes the notation used
throughout this document. The preprint cites three companion
preprints from the same release, \cite{OAI-Companion1}, \cite{OAI-Companion2} and
\cite{OAI-Companion3}, which it cites as \texttt{OpenAIPerfectTwoToOne2026},
\texttt{OpenAIUniformSparsest2026} and \texttt{OpenAIGaussianPropeller2026}; the
first is named as ``not an input to this proof'', and the other two only for
results recorded in its final section. They belong to the same release, so they
are not treated as earlier work here.

\paragraph{What this document reports.}
Each finding is a citation the preprint is missing: part of the preprint's
mathematics (a technique, construction, estimate, proof route or result) that
an identifiable earlier paper already contains, where the preprint does not
cite that paper for it where it is used. ``Earlier'' means publicly available
before the preprint's date. The one exception, overstated novelty, is
described below. Not reported: textbook material in its textbook form; weaker versions of
intermediate statements; and, for content the preprint uses, cases where it
cites a later source that covers that content rather than the original. An
earlier result on the same problem is reported even when later work that the
preprint cites implies it. Findings are ordered by centrality, then by relationship in
the order listed below, then by position in the preprint.

\paragraph{Relationships.}
A finding's relationship is the one between the preprint and the first
earlier work listed for it (the principal), the work with the highest relationship in the
order below; when several hold that relationship, however closely each matches,
the one the others credit or build on, otherwise the earliest. Any other works listed relate to the same content more weakly or
in part; the finding says how, including which work originated the idea when
that is not the first.
\begin{description}[leftmargin=1.5em,font=\normalfont\itshape]
\item[same result.] An earlier paper, or several earlier papers together, prove essentially the same statement.
\item[same technique.] The earlier technique, construction or theory is
reused essentially intact, possibly in a new setting or with new ingredients
around it.
\item[adapted.] The earlier technique's structure is followed, with a change
at its core: a different object, or a key step replaced.
\item[proof shortcut.] A published result would replace or substantially
shorten a large part of the argument.
\item[same route.] An earlier paper proposed the proof strategy, even
heuristically, and the preprint carries it out.
\item[fragment.] One identifiable piece of the earlier technique is reused.
\item[weaker result.] An earlier paper proves a weaker, special-case or
partial-range version of one of the preprint's main results, or of a
statement it presents as a significant contribution.
\end{description}

\paragraph{Centrality.}
Each finding also gets a best-effort judgment of how much of the preprint's
main argument depends on it: \emph{main result} (reserved for earlier work that
covers the main result itself, or most of it), \emph{central},
\emph{supporting} or \emph{peripheral}. An earlier weaker result is never
labelled main result. Each finding states the chain of results it feeds, so
the basis for the judgment is visible.

\paragraph{Citations.}
The ``Cited?'' column of the summary names the earlier work and says what the
preprint cites for the content in question: \emph{not cited} (the work is not
in the bibliography, and nothing the preprint cites credits it for this
content), \emph{cited for other content}, \emph{credited elsewhere} (credited
for this content, directly or through a cited location that credits it, but
not where it is used, in a way that an expert reader who doesn't know the earlier work might not connect with it), or \emph{credited in part}. A citation counts only when it is \emph{for that specific content};
citing the same paper nearby for a different result does not count. Content
the preprint credits where it is used, directly or through a cited location
that itself credits the earlier work (such as a companion
lemma), is not reported, except when its novelty is overstated (below). For
overstated novelty, the column says where the preprint does cite the work.

\paragraph{Flags.}
\begin{description}[leftmargin=1.5em,font=\normalfont\bfseries]
\item[Attribution gap.] Marks the most serious missing citations: the
relationship is same technique, adapted or same result; the centrality is main
result, central or supporting; and the preprint cites the principal earlier
work for this content nowhere.
\item[Overstated novelty.] The preprint does cite the earlier work for this
content, but its abstract, a statement of its significant contributions, or the framing of
a main result presents that content as new, explicitly or in the way a reader
would reasonably take it. This is the one case in which content the preprint
credits is reported.
\end{description}

\paragraph{How to check a claim.}
Locations in the preprint and its companions use its section,
equation and theorem numbers and the page numbers printed in its PDF.
Locations in other sources give the theorem, equation and printed page of the
version listed in the bibliography. Each finding ends with a short checklist
of the passages to put side by side.

\paragraph{Provenance.}
This document was generated by LLM agents (Claude Fable 5.1) running the
precedence-review skill
(\href{https://github.com/jnshi/math-lineage/tree/3499d47/.agents/skills/precedence-review}{jnshi/math-lineage},
commit \texttt{3499d47}), makes no claim of completeness, and has not been
vetted by a human. It was run by JSS.
\space%
\section{The preprint in brief, and notation}\label{sec:overview}

\paragraph{Main result.}
A Unique Games instance has a finite vertex set, a finite alphabet $K$ and a
list of oriented constraints $e=(u_e,v_e,\pi_e)$ with $\pi_e$ a permutation of
$K$; a labeling $a$ satisfies $e$ when $a(v_e)=\pi_e(a(u_e))$, and $\operatorname{val}(G)$ is
the largest fraction (or weight) of satisfied constraints. The preprint's
\ploc{Theorem~1.1 (p.~2)}{00-introduction.tex:22--39} states that for every
fixed $\varepsilon,\delta\in(0,1/2)$ there are an integer $s$ and a
deterministic polynomial-time reduction from 3SAT to explicit unweighted simple
bipartite instances over $K=\mathbb F_2^{s}$ whose constraints are translations
$a(v_e)=a(u_e)+c_e$, with $\operatorname{val}\ge1-\varepsilon$ on satisfiable formulas and
$\operatorname{val}\le\delta$ on unsatisfiable ones. The alphabet and the degree of the
running-time polynomial depend only on $\varepsilon,\delta$.

\paragraph{The route, as the preprint names it.}
The starting hardness input is H\aa stad's nearly satisfiable gap for systems of
3-variable parity equations over $\mathbb F_2$, in weighted form
(\ploc{Theorem~4.1, p.~18}{03-game.tex:25--37}). The construction has five
stages.
\begin{enumerate}
\item \emph{The latent alphabet gadget} (\ploc{\S3, pp.~7--16}{02-latent.tex:1--760}).
  \ploc{Lemma~3.1 (p.~7)}{02-latent.tex:8--31} gives, for every $p_*>0$, a rank
  threshold $r_*$ and, for every large $\ell$, a binary space $\mathcal V\supseteq
  K=\mathbb F_2^{\ell}$, a map $C:\mathcal V\to K$ with $C(x+k)=C(x)+k$, and a noise law
  $\mu$ on $\mathcal V$ with $\Pr[C(x+a)\ne C(x)]\le p_*$ for $a\sim\mu$ and
  uniform $x$, while every linear family $G\le\mathcal V^*$ of restriction rank
  $\dim(G|_K)\ge r_*$ detects $a$ with probability at least $1/8$. The gadget is
  built from quadratic blocks over a finite field of characteristic two,
  concatenated recursively.
\item \emph{The outer questions and the permutation test}
  (\ploc{\S4, pp.~17--21}{03-game.tex:1--361}). An equation question
  $U=(e_1,\dots,e_k)$ is a $k$-tuple of parity equations sampled with
  replacement; a projected question $O$ keeps each equation with probability
  $1-\beta$ and otherwise one of its variables. A \emph{table} on a question $Q$
  is an affine map from its answer space to $\mathcal V$, written as a linear map
  $P\in\operatorname{Hom}(L_Q,\mathcal V)$ with a homogeneous coordinate. Tables are identified
  by an exact \emph{key} (their retained coordinates and induced affine map),
  and the $K$-translates of a key form one vertex (folding); $F_U(P)$ denotes the
  restored answer of a vertex labeling at the table $P$. The instance
  $\mathcal G$ samples $U$, a uniform table $P$, $a\sim\mu$ and a uniform
  $l\in\mathbb F_2^{m}$ and checks $F_U(P)=F_U(P+al^{\top})$
  (\ploc{(4.6), p.~21}{03-game.tex:302--305}); each sample is a translation
  constraint between two vertices.
\item \emph{Matrix extraction and Fourier decoding} (\ploc{\S5, pp.~22--32}{04-matrix.tex:1--767}).
  The structural input is \ploc{Theorem~5.1 (p.~22)}{04-matrix.tex:33--60}, an
  inverse theorem for the degree-two matrix shortcode: a function on
  $\ell\times m$ binary matrices that often keeps its value under a uniform
  rank-one perturbation agrees with an affine evaluation $M\mapsto Mz+u$ on a
  \emph{slice} cut out by at most $r_{\rm s}$ row and $r_{\rm s}$ column
  equations. A spectral comparison (\ploc{Lemma~5.4, p.~25}{04-matrix.tex:244--252})
  turns acceptance of the latent test into acceptance of the uniform rank-one
  test in $K$. The \emph{advice experiment}
  (\ploc{Definition~5.2, p.~23}{04-matrix.tex:168--187}) sends both provers a
  uniform map $A:K\to\mathbb F_2^{r_{\rm s}}$ and the rows $AM$ of the $K$-component of
  a table on $O$, pulled back to $U$; an erasure argument and a comparison of
  the pulled-back table law with the unrestricted law (with $\beta=k^{-2/3}$)
  transfer the slice from $U$ to $O$, and a Fourier decoder on the $O$-prover's
  row fiber yields answer agreement with probability at least a fixed
  $\gamma>0$, independently of $k$.
\item \emph{Clean coordinates and the soundness gap} (\ploc{\S6, pp.~33--35}{05-clean.tex:1--282}).
  On a NO instance the same advice experiment has agreement at most
  $\exp(-ck^{1/3})$: singleton coordinates on which every advice row has zero
  coefficient are \emph{clean}, and after conditioning on the remaining data
  the clean coordinates form an ordinary product of equation-versus-variable
  games, bounded by the projection-game repetition theorem of Dinur and
  Steurer (\ploc{Theorem~6.1, p.~33}{05-clean.tex:26--42}). Together with
  stage~3 this gives a matrix game with value $\le 99/100$ on NO instances and
  completeness $1-k\xi-p_*/2$ on YES instances
  (\ploc{Proposition~6.3, p.~36}{05-clean.tex:243--256}).
\item \emph{Completion} (\ploc{\S7, pp.~36--38}{06-finish.tex:1--241}).
  Weights are rounded and every edge subdivided into four
  (\ploc{Lemma~7.1, p.~36}{06-finish.tex:17--30}), giving an unweighted simple
  bipartite game with a fixed, alphabet-independent gap, and one alphabet-independent
  parallel repetition brings the soundness below $\delta$.
\end{enumerate}
\ploc{\S8 (pp.~39--42)}{07-consequences.tex:1--331} and
\ploc{Appendix~A (pp.~43--53)}{finite-consequences.tex:1--988} derive
approximation consequences from Theorem~1.1 through earlier conditional
reductions, which the preprint attributes to their authors.

\paragraph{Notation used in this document.}
$\mathbb F_2$ is the binary field; $\operatorname{Hom}(E,F)$ the linear maps $E\to F$, and
$E^*=\operatorname{Hom}(E,\mathbb F_2)$. $K=\mathbb F_2^{\ell}$ is the Unique Games alphabet,
$\mathcal V=K\oplus W$ the latent space with the fixed complement $W$, $C$ and
$\mu$ the latent map and noise law, $p_*$ the change probability and $r_*$ the
rank threshold of Lemma~3.1. $k$ is the tuple length, $\beta$ the thinning
probability, $U$ an equation question and $O$ its projection, with
$\pi$ (answers) and $\bar\pi$ (linear extensions) the projections from $U$ to
$O$. $P=(M,T)$ splits a table into its $K$- and $W$-valued components;
$r_{\rm s}$ is the slice parameter of Theorem~5.1; $A$ is the advice map and
$\{M':AM'=AM\}$ a row fiber. Fourier coefficients are normalized,
$\widehat f(g)=\E_x f(x)(-1)^{g(x)}$. Companions are cited as the preprint
cites them; their locators use their own numbering.
\space%

% IDs are assigned by LaTeX in the order of the \input lines (F1, F2, ...).
% Order by magnitude: centrality (main result, central, supporting,
% peripheral), then the relationship (same result, same technique, adapted,
% proof shortcut, same route, fragment, weaker result), then reading order.
% precedence_stats.py checks this order.
% Withdrawing or reordering a finding is just deleting or moving its \input line.
\section{Findings}\label{sec:findings}
\findinggroup{F}
\finding[Erasure argument for the covering fraction]{Re-randomizing a function on its structured pieces to show that random advice hits one}{erasure-covering-argument}

\paragraph{Summary.}
The covering half of \ploc{Lemma~5.5 (pp.~26--27)}{04-matrix.tex:325--415} follows the ``erasure'' argument of Dinur, Khot, Kindler, Minzer and Safra, \sloc{DinurEtAl2018-Towards2to1}{Lemma~3.7 with Claims~3.8--3.9, pp.~23--24}, already in their December 2016 ECCC version \sloc{DinurEtAl2016-Towards2to1-ECCC}{Lemma~3.7, pp.~14--16}. The preprint cites that paper only in the survey of \ploc{\S1.3 (p.~3)}{00-introduction.tex:119--122}, for its conditional 2-to-1 reduction, and nowhere in \ploc{\S5.4}{04-matrix.tex:299--443}, although it calls the step an ``erasure argument'' (\ploc{\S5.4, p.~26}{04-matrix.tex:305}: ``The erasure argument below''; \ploc{\S1.4, p.~5}{00-introduction.tex:248--249}: ``An erasure argument''), the earlier paper's own word for its device. In both, an inverse statement (Hypothesis~3.6 there, \ploc{Theorem~5.1, p.~22}{04-matrix.tex:33--54} here) yields one structured piece, while decoding with shared advice needs \emph{random} advice to hit one with positive probability, the mass \ploc{(5.8)}{04-matrix.tex:329--335} behind \ploc{Proposition~5.3}{04-matrix.tex:197}. The insight, due to Dinur et al., is to re-randomize the function on the structured pieces: random values cannot create agreement with any structured target (concentration plus a union bound), a small erased set barely moves the test acceptance, so the inverse statement forces the pieces to cover a constant fraction of the domain. \rel{same technique}: every ingredient recurs, down to the expectation $2^{-\ell}$ of a random match; the preprint erases all good pieces at once and argues by contradiction where the earlier proof erases one piece at a time.

\paragraph{Correspondence.}
$f=F_U(\cdot,T)$; a row fiber is \emph{good} when it contains a slice $\mathcal S$, of size $n_{\mathcal S}$, on which $f$ agrees with some $Mz+u$ with probability at least $\alpha/2$ (\ploc{p.~26}{04-matrix.tex:308--316}); $B$ is the union of the good fibers; $\eta=2^{-r_*}/16$. In the earlier paper $F[\cdot]$ assigns a linear function to each $\ell$-dimensional $L\subseteq\mathbb F_2^{n}$, and Hypothesis~3.6 \sloc{DinurEtAl2018-Towards2to1}{(3.1), p.~22} turns $\operatorname{agreement}(F)\ge\delta$ into a zoom pair $Q\subseteq W$ and a linear $g_{Q,W}$ on $W$ with $\Pr_L[g_{Q,W}|_L=F[L]\mid Q\subseteq L\subseteq W]\ge C$.
\begin{center}\small
\begin{xltabular}{\linewidth}{@{}>{\raggedright\arraybackslash}p{0.12\linewidth}>{\raggedright\arraybackslash}X>{\raggedright\arraybackslash}X@{}}
\toprule
 & Preprint & Dinur et al.\ \cite{DinurEtAl2018-Towards2to1}\\
\midrule
\endfirsthead
\toprule
 & Preprint & Dinur et al.\ \cite{DinurEtAl2018-Towards2to1}\\
\midrule
\endhead
Erasure & replace $f$ at the entries of $B$ by independent uniform elements of $K$ \ploc{p.~26}{04-matrix.tex:358--362} & a random linear $F[L]$ on every $L\supseteq Q$ for each $Q$ found (Lemma~3.7, p.~23)\\
\addlinespace
Acceptance barely moves & acceptance moves by at most $2\Pr[B]<2\eta$, from $\ge4\eta$ to $\ge2\eta$ \ploc{p.~26}{04-matrix.tex:360--365} & after $N$ pieces a test pair meets an erased $L$ with probability $\le2N\cdot2^{\ell^{2}-qn}$, so agreement stays $\ge\delta/2$ while $N$ is small (Claim~3.8, pp.~23--24)\\
\addlinespace
Random values create no structure & matches at randomized entries: mean $\le2^{-\ell}n_{\mathcal S}$, Hoeffding $\exp(-\alpha^{2}n_{\mathcal S}/32)$, union bound over $2^{(2r+1)(\ell+m)}$ descriptions \ploc{p.~27}{04-matrix.tex:367--394} & a match on an erased $L$ has $\E Z_L=2^{-\ell}$; Chernoff; union bound over $2^{qn}2^{rn}2^{n-r}$ triples $(Q,W,g_{Q,W})$ (Claim~3.9, p.~24)\\
\addlinespace
Covering conclusion & agreement $<\alpha$ contradicts Theorem~5.1, so $\Pr[B]\ge\eta$; a uniform $A$ hits a good fiber with probability $2^{-\ell r_{\rm s}}$, whence (5.8) \ploc{p.~27}{04-matrix.tex:396--414} & each piece keeps $C-2^{1-\ell}\ge C/2$; at the end at least an $\alpha(\ell)=\delta2^{-\ell^{2}-2}$ fraction of the $Q$ are covered (pp.~23--24)\\
\addlinespace
\bottomrule
\end{xltabular}
\end{center}

\paragraph{How to check.}
\begin{enumerate}
\item The proof opening of \sloc{DinurEtAl2018-Towards2to1}{Lemma~3.7, p.~23} (random reassignment on every $L\supseteq Q$) against \ploc{p.~26}{04-matrix.tex:358--362} (uniform replacement on $B$).
\item \sloc{DinurEtAl2018-Towards2to1}{Claim~3.8, pp.~23--24} (reassigned pairs: probability $\le2N\cdot2^{\ell^{2}-qn}$) against \ploc{p.~26}{04-matrix.tex:363--365} (acceptance moves by $\le2\Pr[B]$).
\item \sloc{DinurEtAl2018-Towards2to1}{Claim~3.9, p.~24} (expectation $2^{-\ell}$, Chernoff, union bound) against \ploc{p.~27}{04-matrix.tex:380--394} (mean $2^{-\ell}n_{\mathcal S}$, Hoeffding, union bound).
\end{enumerate}

\space%
\finding[Shared zoom-in, private zoom-out, random extension]{Decoding with a shared zoom-in sent as advice, a zoom-out each prover finds in its own table, and a uniformly random extension as the answer}{zoom-in-advice-zoom-out-extension}

\paragraph{Summary.}
Dinur, Khot, Kindler, Minzer and Safra (DKKMS)
\sloc{DinurEtAl2018-Towards2to1}{\S1.4, p.~13; \S5, p.~41} decode with a prover
rule of this shape: the verifier sends the zoom-in part of the structure as
shared advice chosen independently of the tables, each prover finds the zoom-out
part in its own table, and the answer is a uniformly random linear extension
over the zoom-out, valid by invariance and costing a factor $2^{-r}$ (the
shared-advice sentence and the list paragraph are already in their December 2016
version \sloc{DinurEtAl2016-Towards2to1-ECCC}{p.~8}).
The preprint cites DKKMS only for their conditional 2-to-1 reduction
(\ploc{\S1.3, p.~3}{00-introduction.tex:119--122}) and credits the advice to
Khot, Minzer and Safra (2017)
\sloc{KhotMinzerSafra2017-IndependentSets2to2Grassmann}{\S\S3.3--3.4, pp.~20--21}
at \ploc{\S1.4, p.~5}{00-introduction.tex:279--282}, \ploc{\S5.5, p.~27}{04-matrix.tex:445--449}
and \ploc{\S6, p.~33}{05-clean.tex:8--9}; the 2017 game has a zoom-in only, and
its provers answer with a global function
(\sloc{KhotMinzerSafra2017-IndependentSets2to2Grassmann}{Hypothesis~2.12, p.~19; \S5.1, p.~29}).
The preprint's $U$-prover (\ploc{\S5.4--\S5.6, pp.~26--32}{04-matrix.tex:299--749})
has the DKKMS shape: the advice $(A,AM)$ fixes the row equations $AM=S_0$, the
zoom-in by the proof of \ploc{Theorem~5.1, p.~22}{04-matrix.tex:114--118}; the
prover chooses the column equations $Mq_j=t_j$, the zoom-out, from its own table;
and it answers with a uniform element of the coset (5.9), of size at most
$2^{r_{\rm s}}$ and valid by folding. \rel{adapted}: the shape is DKKMS's, but the
preprint drops what DKKMS call their main technical contribution, lists of
maximal zoom-outs with a size bound and a common-element argument
(\sloc{DinurEtAl2018-Towards2to1}{p.~13; Lemma~3.17, pp.~30--31}): the $U$-prover
picks one witness by a fixed order, and the $O$-prover forms no list but samples
squared Fourier mass on its own row fiber after H{\aa}stad and Khot--Safra. Credit
status: credited in part (the shared zoom-in advice, to the 2017 paper); DKKMS
cited for other content.

\paragraph{Correspondence.}
\begin{center}\small
\begin{xltabular}{\linewidth}{@{}>{\raggedright\arraybackslash}p{0.13\linewidth}>{\raggedright\arraybackslash}X>{\raggedright\arraybackslash}X@{}}
\toprule
 & Preprint & DKKMS \cite{DinurEtAl2018-Towards2to1}\\
\midrule
\endfirsthead
\toprule
 & Preprint & DKKMS \cite{DinurEtAl2018-Towards2to1}\\
\midrule
\endhead
Shared zoom-in & uniform $A$ and rows $S_O=AM_O$, sent to both provers independently of the tables (\ploc{Definition~5.2, (5.5), p.~24}{04-matrix.tex:168--187}) & $Q$ chosen by the verifier and sent to both provers as ``shared advice'' (p.~13; \S5, p.~41)\\
\addlinespace
Private zoom-out & column equations $Mq_j=t_j$ and target chosen ``using a fixed order'' from the row advice alone (\ploc{Lemma~5.5, pp.~26--27}{04-matrix.tex:339--354,438--440}) & $(f,W)$, $W\supseteq Q$ of codimension $\le r$, chosen from $\operatorname{LIST}_Q(F_U)$ ``from her/his own viewpoint'' (pp.~13, 41)\\
\addlinespace
Answer and loss & uniform $a_U\in\mathcal A=(z+Z)\cap\{e_0^\top v=1\}$, $|\mathcal A|\le2^{r_{\rm s}}$; success probability $\ge1/|\mathcal A|\ge2^{-r_{\rm s}}$ (\ploc{Equation~(5.9), p.~26; \S5.6, p.~30; Lemma~5.8, p.~32}{04-matrix.tex:345--350,632--636,745--747}) & ``chooses $r$ bits to randomly extend $f$ to a linear function on all of $X_U$'' (p.~41); win ``with probability at least $2^{-i}\ge2^{-r}$'' (p.~43)\\
\addlinespace
\bottomrule
\end{xltabular}
\end{center}

\paragraph{How to check.}
\begin{enumerate}
\item The paragraph ``To handle the zoom-outs'' at \sloc{DinurEtAl2018-Towards2to1}{p.~13}
and ``Strategies of the players'' at \sloc{DinurEtAl2018-Towards2to1}{\S5, p.~41}
against the $U$-prover's rule at \ploc{\S5.6, p.~30}{04-matrix.tex:632--636} and
\ploc{Equation~(5.9), p.~26}{04-matrix.tex:345--350}.
\item \sloc{DinurEtAl2018-Towards2to1}{pp.~42--43} (``regardless of the $i\le r$ coin
tosses''; ``$2^{-i}\ge2^{-r}$'') against the coset step at \ploc{p.~27}{04-matrix.tex:416--442}
and the factor $1/|\mathcal A|\ge2^{-r_{\rm s}}$ at \ploc{Lemma~5.8, p.~32}{04-matrix.tex:745--747}.
\item The credits at \ploc{p.~5}{00-introduction.tex:279--282}, \ploc{p.~27}{04-matrix.tex:445--449}
and \ploc{p.~33}{05-clean.tex:8--9} against
\sloc{KhotMinzerSafra2017-IndependentSets2to2Grassmann}{\S\S3.3--3.4, pp.~20--21; \S5.1, p.~29}:
an advice game whose prover outputs $g_Q$, with no zoom-out.
\end{enumerate}

\space%
\finding[Folded vertex classes with a representative and an unfolding rule]{A folding written out as a constraint graph: vertex classes, one representative per class, an unfolding rule and aggregated weighted constraints}{folded-vertex-classes}

\paragraph{Summary.}
The nearest written-out precedent for the vertex set of the preprint's permutation
instance (\ploc{\S\S4.3--4.4, pp.~20--21}{03-game.tex:165--317}) is the folded
constraint graph $G_{\rm folded}$ of Dinur, Khot, Kindler, Minzer and Safra
(DKKMS) \sloc{DinurEtAl2018-Towards2to1}{\S1.3.3, pp.~9--10, and \S4.2, pp.~33--35},
first circulated in December 2016
\sloc{DinurEtAl2016-Towards2to1-ECCC}{\S1.3.3, p.~6, and \S4.2, p.~24}, for the
same $\beta$-thinned $k$-tuple outer game. The preprint cites DKKMS once, for
their reduction ``conditional on an agreement hypothesis''
(\ploc{\S1.3, p.~3}{00-introduction.tex:119--122}), and not there; at
\S\S4.3--4.4 it credits Khot and Safra
\sloc{KhotSafra2013-TwoProverStrongSoundness}{\S1.1, pp.~867--868} for
identifying table positions across questions and Bellare, Goldreich and Sudan
\sloc{BellareGoldreichSudan1998-FreeBitsPCPs}{\S3.3, pp.~31--32} for folding a
table over translations with a canonical representative, which supply its
skeleton: two tables are one vertex exactly when their exact keys are equal,
and a $K$-orbit is represented by its key with zero $K$-component of the
intercept. Shared with DKKMS beyond those principles is the packaging of a
folding as a constraint graph: classes of question--table pairs whose honest
answers determine one another, a label read at one representative and unfolded
to every member, and the unfolded test's constraints pushed forward to the
classes with their weights summed. \rel{fragment}: only this
generic packaging is shared; DKKMS's own machinery (their subspace relation,
Lemma~4.2's arbitrary representative $R$, Claim~4.4) has no counterpart in the
preprint.

\paragraph{Correspondence.}
\begin{center}\small
\begin{xltabular}{\linewidth}{@{}>{\raggedright\arraybackslash}p{0.14\linewidth}>{\raggedright\arraybackslash}X>{\raggedright\arraybackslash}X@{}}
\toprule
 & Preprint & DKKMS \cite{DinurEtAl2018-Towards2to1}\\
\midrule
\endfirsthead
\toprule
 & Preprint & DKKMS \cite{DinurEtAl2018-Towards2to1}\\
\midrule
\endhead
Class and representative & the tables $(Q,P)$ with one exact key, with their $K$-translates $P+he_0^{\top}$; represented by the unique key in the orbit with zero $K$-component of the intercept \ploc{pp.~20--21}{03-game.tex:233--267} & $\mathcal C(U_0,L_0)=\{(U,L)\in A:\ L+H_U+H_{U_0}=L_0+H_U+H_{U_0}\}$, $H_U$ the span of the equations of $U$; represented by an $\ell$-dimensional $R$ with $L+H_U=R+H_U$ for all members, fixed arbitrarily (\S4.2 and Lemma~4.2, p.~34)\\
\addlinespace
Unfolding & $F_Q(P)=\lambda(v(Q,P))+c(Q,P)$, $\lambda$ a labeling of the orbit vertices and $c(Q,P)$ the $K$-component of the intercept; hence $F_Q(P+he_0^{\top})=F_Q(P)+h$ \ploc{(4.4), p.~21}{03-game.tex:273--282} & $\mathcal A(U,L)=\tilde\sigma_U|_L$, $\tilde\sigma_U(x+y)=\sigma(x)+h_U(y)$ for $x\in R$, $y\in H_U$, with $h_U(x_e)=b_e$ the right-hand sides (Definition~4.3, p.~35)\\
\addlinespace
Aggregated constraints & a sample $(U,P,a,l)$ is an edge between the orbit vertices of $P$ and $P+al^{\top}$ with constraint $\lambda(v_1)=\lambda(v_0)+c_0+c_1$, weighted by its sampling probability \ploc{\S4.4, p.~21}{03-game.tex:307--314} & $\tilde E$: sample $((U,L),(V,L'))\in E$, output $(\mathcal C(U,L),(V,L'))$; constraints ``the aggregate of the constraints in $G_{\rm unfolded}$ \dots\ The weights are summed up as well'' (pp.~34--35)\\
\addlinespace
\bottomrule
\end{xltabular}
\end{center}

\paragraph{How to check.}
\begin{enumerate}
\item The key, ``two tables have the same key exactly when \dots''
\ploc{\S4.3, p.~20}{03-game.tex:233--237}, and the orbit representative
\ploc{p.~21}{03-game.tex:261--267} against $\mathcal C(U_0,L_0)$ and
\sloc{DinurEtAl2018-Towards2to1}{Lemma~4.2, p.~34}.
\item \ploc{(4.4), p.~21}{03-game.tex:273--282} against
\sloc{DinurEtAl2018-Towards2to1}{Definition~4.3, p.~35}.
\item The edges and weights of $\mathcal G$, \ploc{\S4.4, p.~21}{03-game.tex:307--314},
against the ``aggregate'' paragraph (p.~35).
\end{enumerate}

\space%
\finding[Identity-target kernel clustering and its rounding constant]{The identity-target kernel clustering problem and the Gaussian rounding constant \texorpdfstring{$(8\pi/9)(1-1/k)$}{(8 pi/9)(1 - 1/k)}}{kernel-clustering-khot-naor}

\paragraph{Summary.}
The maximization problem defined in the preprint's ``Kernel clustering''
paragraph \ploc{\S8.3, p.~43}{07-consequences.tex:306--331} and the constant
$\alpha_k=\frac{8\pi}{9}(1-\frac1k)$ it calls ``the propeller rounding factor''
are Khot and Naor's \cite{KhotNaor2008-ApproximateKernelClustering}, whom the
preprint cites nowhere: the paragraph's only citation is the companion's
\ploc{Theorem~6.1 (p.~18)}{kernel-application.tex:53--64}, and the sentence
opening the preprint's \S8, ``The underlying conditional reductions and
approximation thresholds are due to the cited authors''
\ploc{\S8, p.~39}{07-consequences.tex:3--6}, names for this paragraph only that
companion. The problem is the kernel clustering problem of Song, Smola, Gretton
and Borgwardt with the identity target matrix, as Khot and Naor set it up
\sloc{KhotNaor2008-ApproximateKernelClustering}{pp.~1--2, eq.~(2)}; $\alpha_k$
is the loss factor of their Gaussian rounding of the semidefinite relaxation
\sloc{KhotNaor2008-ApproximateKernelClustering}{Theorem~2.1, p.~8}, which they
conjectured to be the sharp Unique Games threshold for every $k\ge3$
\sloc{KhotNaor2008-ApproximateKernelClustering}{p.~2}, having shown that this
threshold is at least $(1-1/k)/C(k)$ for a Gaussian partition constant $C(k)$
\sloc{KhotNaor2008-ApproximateKernelClustering}{pp.~5 and~7; \S3.3, pp.~26--28},
so that $\alpha_k$ is sharp exactly when $C(k)=9/(8\pi)$, as they proved for
$k=3$. The companion credits Khot and Naor for the problem, the constant and the
reduction (\ploc{\S1.3, p.~3}{main.tex:151--167};
\ploc{\S6.1, pp.~17--18}{kernel-application.tex:5--51}), but it belongs to the
same release, so a reader of the preprint is not pointed to their 2008 paper
(journal version 2009). \rel{same result}: the preprint states Khot and Naor's
problem and their rounding constant; the hardness statement it records is the
companion's and is cited.

\paragraph{Correspondence.}
\begin{center}\small
\begin{xltabular}{\linewidth}{@{}>{\raggedright\arraybackslash}p{0.13\linewidth}>{\raggedright\arraybackslash}X>{\raggedright\arraybackslash}X@{}}
\toprule
 & Preprint & Khot--Naor \cite{KhotNaor2008-ApproximateKernelClustering}\\
\midrule
\endfirsthead
\toprule
 & Preprint & Khot--Naor \cite{KhotNaor2008-ApproximateKernelClustering}\\
\midrule
\endhead
Problem & rational symmetric $A\succeq0$ with $A\one=0$; $\operatorname{OPT}_k(A)=\max_\sigma\sum_{p,q}a_{pq}\one_{\{\sigma(p)=\sigma(q)\}}$ over $\sigma:\{1,\dots,N\}\to\{1,\dots,k\}$, $k\ge3$ fixed \ploc{(p.~43)}{07-consequences.tex:308--320} & centered kernel $A\succeq0$, $\sum_{i,j}a_{ij}=0$, target $B=I_k$: $\operatorname{Clust}(A|I_k)=\max_\sigma\sum_{i,j=1}^n a_{ij}\one_{\{\sigma(i)=\sigma(j)\}}$ (eq.~(2), p.~2; \S2, p.~7)\\
\addlinespace
Loss factor & $\alpha\ge1$: returned value $\ge\operatorname{OPT}_k(A)/\alpha$ \ploc{(p.~43)}{07-consequences.tex:320--321} & $\operatorname{Clust}(A|B)\le\frac{8\pi}{9}(1-\frac1k)\,\E\bigl[\sum_{i,j}a_{ij}b_{\lambda(i)\lambda(j)}\bigr]$ for the rounding $\lambda$; ``The same bound in (12) holds true if $B$ is the identity matrix'' (Theorem~2.1, eq.~(12), p.~8)\\
\addlinespace
Constant & $\alpha_k=\frac{8\pi}{9}(1-\frac1k)$, ``the propeller rounding factor'' \ploc{(p.~43)}{07-consequences.tex:324--331} & $\frac{8\pi}{9}(1-\frac1k)$: approximation factor for $\operatorname{Clust}(A|I_k)$, conjectured sharp UGC threshold for $k\ge3$ (p.~2; Theorem~1.1, p.~4)\\
\addlinespace
\bottomrule
\end{xltabular}
\end{center}

\paragraph{How to check.}
\begin{enumerate}
\item The problem definition \ploc{(p.~43)}{07-consequences.tex:306--321} against
\sloc{KhotNaor2008-ApproximateKernelClustering}{p.~1 (centered kernel) and eq.~(2), p.~2}
with $B=I_k$.
\item The constant $\alpha_k$ \ploc{(p.~43)}{07-consequences.tex:324--331} against
\sloc{KhotNaor2008-ApproximateKernelClustering}{eq.~(12) and the sentence after it, p.~8}
and the approximation-factor and conjecture sentences on
\sloc{KhotNaor2008-ApproximateKernelClustering}{p.~2}.
\item The citation list of the paragraph (\ploc{p.~43}{07-consequences.tex:329}: the
companion's Theorem~6.1 only) against the opening sentence
\ploc{\S8, p.~39}{07-consequences.tex:3--6} and
the companion's credits \ploc{\S1.3, p.~3}{main.tex:163--167} and
\ploc{\S6.1, pp.~17--18}{kernel-application.tex:5--6, 40--51}.
\end{enumerate}

\space%
\finding[Correlation Clustering from Multicut hardness]{Hardness of weighted Correlation Clustering inherited from Multicut hardness through the known equivalence of the two problems}{correlation-clustering-via-multicut}

\paragraph{Summary.}
Chawla, Krauthgamer, Kumar, Rabani and Sivakumar, whose Multicut hardness the
preprint's \ploc{Corollary~8.3 (p.~42)}{07-consequences.tex:261--296} rests on,
already state that their results ``immediately extend also to the Correlation
Clustering problem \ldots\ of minimizing disagreements in a weighted graph,
since the approximability of this problem is known to be equivalent to within
constant factors to that of Multicut in weighted graphs''
\sloc{ChawlaEtAl2006-HardnessMulticutSparsestCut}{p.~2}, crediting the
equivalence to Charikar, Guruswami and Wirth and to Emanuel and Fiat. The
preprint cites or points to that paper five times in
\ploc{\S8.3}{07-consequences.tex:234--296} (once before the corollary, four
times in the proof), but never for this extension: it writes that
Corollary~1.3 of that paper ``gives the first three conclusions'' and derives
the fourth, Correlation Clustering, as its own further step
(\ploc{proof of Corollary~8.3, pp.~42--43}{07-consequences.tex:272--296}),
crediting only the transformation it uses, to Demaine, Emanuel, Fiat and
Immorlica \sloc{DemaineEtAl2006-CorrelationClusteringGeneralWeightedGraphs}{Thm~4.7
and Cor.~4.8, p.~14}, the journal version of the Emanuel--Fiat paper
(\sloc{DemaineEtAl2006-CorrelationClusteringGeneralWeightedGraphs}{p.~1}).
The shared approach is to prove nothing new for Correlation Clustering: a
Multicut instance becomes a Correlation Clustering instance once its edges are
labelled $+$ and each demand pair gets a $-$ edge heavier than all $+$ edges
together, so constant-factor hardness transfers; the equivalence is
Charikar--Guruswami--Wirth's and Emanuel--Fiat's, and applying it to the Unique
Games Multicut instances is the Multicut source's. \rel{same route}: the
source proposed this step for these very instances, and the preprint carries
it out with the credited transformation plus one check of its own, that the
demand pairs are nonedges of the Multicut graph
(\ploc{p.~42}{07-consequences.tex:285--290}), which neither earlier paper states.

\paragraph{Correspondence.}
\begin{center}\small
\begin{xltabular}{\linewidth}{@{}>{\raggedright\arraybackslash}p{0.17\linewidth}>{\raggedright\arraybackslash}X>{\raggedright\arraybackslash}X@{}}
\toprule
 & Preprint & Chawla et al.\ \cite{ChawlaEtAl2006-HardnessMulticutSparsestCut}; Demaine et al.\ \cite{DemaineEtAl2006-CorrelationClusteringGeneralWeightedGraphs}\\
\midrule
\endfirsthead
\toprule
 & Preprint & Chawla et al.\ \cite{ChawlaEtAl2006-HardnessMulticutSparsestCut}; Demaine et al.\ \cite{DemaineEtAl2006-CorrelationClusteringGeneralWeightedGraphs}\\
\midrule
\endhead
Cut hardness used & ``Corollary~1.3 of [the source] gives the first three conclusions'' \ploc{(p.~42)}{07-consequences.tex:272--273} & Cor.~1.3, p.~4: Multicut, Sparsest-Cut, Min-$2\mathrm{CNF}^{\equiv}$ Deletion within every constant $L$ under the conjecture\\
\addlinespace
Passage to Correlation Clustering & by the Demaine et al.\ transformation, ``so the same fixed-factor hardness transfers'' \ploc{(pp.~42--43)}{07-consequences.tex:285--296} & ``Our results immediately extend also to the Correlation Clustering problem \ldots\ equivalent to within constant factors to that of Multicut'' (Chawla et al.\ p.~2)\\
\addlinespace
The transformation & ``preserves the optimum and converts every clustering into a multicut of no greater cost'' \ploc{(p.~43)}{07-consequences.tex:293--295} & $G_H=H$ with $+$ labels plus a $-$ edge $(s_i,t_i)$ of weight $\sum_{e\in H}w(e)+1$ per pair $\langle s_i,t_i\rangle$ (Demaine et al.\ \S4.2.1, p.~13); Thm~4.7 and Cor.~4.8 (p.~14)\\
\addlinespace
\bottomrule
\end{xltabular}
\end{center}

\paragraph{How to check.}
\begin{enumerate}
\item The sentence ``Our results immediately extend also to the Correlation
Clustering problem\ldots'' at
\sloc{ChawlaEtAl2006-HardnessMulticutSparsestCut}{p.~2} against ``gives the
first three conclusions'' and the Correlation Clustering paragraph of the
\ploc{proof of Corollary~8.3, pp.~42--43}{07-consequences.tex:272--296}.
\item The construction and Theorem~4.7 at
\sloc{DemaineEtAl2006-CorrelationClusteringGeneralWeightedGraphs}{\S4.2.1,
pp.~13--14} against ``preserves the optimum and converts every clustering
into a multicut of no greater cost'' (\ploc{p.~43}{07-consequences.tex:293--295}).
\end{enumerate}

\space%
\finding[Earlier Unique Games hardness points]{Earlier NP-hardness points, and conditional proofs, for the Unique Games problem}{earlier-unique-games-hardness-points}

\paragraph{Summary.}
Before the preprint, NP-hardness of the Unique Games problem that
\ploc{Theorem~1.1 (p.~2)}{00-introduction.tex:22--39} settles had been
proved unconditionally at several weaker completeness--soundness points,
and the Unique Games Conjecture had been derived from other hypotheses:
O'Donnell and Wright prove hardness of $1/2$-satisfiable versus
$(3/8+\epsilon)$-satisfiable instances
\sloc{ODonnellWright2012-NewPointNPHardnessUniqueGames}{abstract, p.~1};
H\aa stad, Huang, Manokaran, O'Donnell and Wright
\cite{HastadEtAl2015-ImprovedNPInapprox2LinearEquations}, Wiman
\cite{NL-Wiman2018} and Martinsson
\cite{Martinsson2024-NPHardnessApproximationCurveMax2Lin2} prove
$1-\epsilon$ versus $1-C\epsilon$ over the binary alphabet; Feige and
Reichman \cite{NL-FeigeReichman2004} prove hardness by a factor
$q^{\beta}$ at low completeness; Bhangale and Khot
\cite{BhangaleKhot2019-UGHardnessToNPHardnessByLosingHalf} prove the
completeness-$1/2$ point in the strong-left, left-regular form the preprint
uses in its \S8 and Appendix~A; and Raghavendra and Steurer
\cite{RaghavendraSteurer2010-GraphExpansionUGC} and Eldan and Moshkovitz
\cite{EldanMoshkovitz2020-NonUniqueToBooleanUniqueGames} derive the
conjecture, or a Boolean form of it, from the Small-Set Expansion and the
Subspaces Near-Intersection hypotheses. The preprint cites none of these
eight works. Its survey of earlier work
(\ploc{\S\S1.2--1.3, pp.~2--3}{00-introduction.tex:57--131}) names a single
earlier hardness point, the completeness-$1/2$ consequence of the 2-to-2
line (\ploc{\S1.3, p.~3}{00-introduction.tex:129--131}); and H\aa stad's
paper, cited for the 3-variable parity gap and its Fourier decoding
(\ploc{\S1.2, p.~2}{00-introduction.tex:72};
\ploc{Theorem~4.1, p.~18}{03-game.tex:25--40};
\ploc{\S5.6, p.~30}{04-matrix.tex:621--623}), contains the long-standing
best earlier binary point, factor $12/11$ for two-variable equations
\sloc{Hastad2001-SomeOptimalInapproximability}{Thm~5.14, p.~35}, which
O'Donnell and Wright call ``the best previous result''
\sloc{ODonnellWright2012-NewPointNPHardnessUniqueGames}{p.~3}.
\rel{weaker result}: each statement is a partial-range, fixed-alphabet or
conditional version of Theorem~1.1, which gives completeness $1-\epsilon$
and soundness $\delta$ for every pair $\epsilon,\delta$. O'Donnell and
Wright are listed first because the later points are measured against
theirs: Dinur, Khot, Kindler, Minzer and Safra present their $1/2$ versus
$1/4+\epsilon$ corollary as improving its soundness
\sloc{DinurEtAl2018-Towards2to1}{Cor.~1.19, p.~15}.

\paragraph{Setup.}
A \emph{point of NP-hardness} $(c,s)$ for a constraint problem means that
distinguishing instances of value at least $c$ from instances of value at
most $s+\epsilon$ is NP-hard for every fixed $\epsilon>0$
\sloc{ODonnellWright2012-NewPointNPHardnessUniqueGames}{Def.~2.3, p.~5}.
Max-2Lin(2), the problem of satisfying equations $x_i+x_j=b$ over
$\mathbb F_2$, is the case $s=1$ of the preprint's translation games
(Section~\ref{sec:overview}); Max-Cut is its subcase $b=1$. The preprint
records that binary translation constraints are an equivalent formulation
of the conjecture (\ploc{\S1.1, p.~2}{00-introduction.tex:44--46}), citing
Khot, Kindler, Mossel and O'Donnell's Corollary~13 and \S11.1; the
equivalence itself is their Corollary~15
\sloc{KhotEtAl2005-OptimalInapproxMaxCut}{Cor.~13 and Cor.~15, p.~15; \S11.1, p.~26}.
The conjecture's point is $(1-\epsilon,\epsilon)$ for every $\epsilon$ with
the alphabet chosen afterwards. None of the results below enters the proof;
they are the earlier state of the problem and belong beside the survey at
\ploc{\S\S1.2--1.3, pp.~2--3}{00-introduction.tex:57--131}. One recurs in
form: Bhangale and Khot's statement is the preprint's ``strengthened
bipartite formulation'' (\ploc{\S8.1, p.~40}{07-consequences.tex:49--54};
\ploc{Lemma~A.1, p.~43}{finite-consequences.tex:22--34}), whose composition
with the ordering reduction is finding~\ref{f:strong-left-ordering-composition}.

\paragraph{The earlier results.}
The table lists each statement in its source's terms and the point it gives
for Theorem~1.1's problem; the first six rows are unconditional, the last
two conditional. In the near-satisfiable rows the gap is $1-\epsilon$
versus $1-C\epsilon$ with a fixed $C$; Khot, Minzer and Safra note that a
gap $(1-\epsilon,\,1-K(\epsilon)\sqrt{\epsilon})$ with $K(\epsilon)$ growing
would already be boosted by parallel repetition to the full conjecture
\sloc{KhotMinzerSafra2018-PseudorandomGrassmannExpansion}{p.~5}.
\begin{center}\small
\begin{xltabular}{\linewidth}{@{}>{\raggedright\arraybackslash}p{0.2\linewidth}>{\raggedright\arraybackslash}X>{\raggedright\arraybackslash}p{0.27\linewidth}@{}}
\toprule
Work & Statement, in the source's terms & Point for Theorem~1.1's problem\\
\midrule
\endfirsthead
\toprule
Work & Statement, in the source's terms & Point for Theorem~1.1's problem\\
\midrule
\endhead
H\aa stad, 2001 \cite{Hastad2001-SomeOptimalInapproximability} &
``For any $\epsilon>0$ it is NP-hard to approximate Max-E2-Lin-2 within a factor $12/11-\epsilon$'' (Thm~5.14, p.~35), by a gadget reduction from Max-E3-Lin-2 (pp.~35--36); recorded as the point $(\tfrac34,\tfrac{11}{16})$ in \sloc{ODonnellWright2012-NewPointNPHardnessUniqueGames}{Thm~2.4, p.~5}. &
Binary alphabet, $(3/4,\,11/16)$. Cited for the 3-variable parity gap and its Fourier decoding (\ploc{Theorem~4.1, p.~18}{03-game.tex:25--40}; \ploc{\S5.6, p.~30}{04-matrix.tex:621--623}).\\
\addlinespace
Feige and Reichman, 2004 \cite{NL-FeigeReichman2004}, as described in \sloc{KhotEtAl2005-OptimalInapproxMaxCut}{\S11.1, p.~27} &
``NP-hard to approximate MAX-2LIN($q$) to within a factor of $1/q^{\beta}$ for some universal constant $\beta>0$; however this hardness is located at a gap of $\epsilon$ vs.\ $\epsilon/q^{\beta}$''; near satisfiability, $1-\eta$ versus $1-C\eta$ ``for some relatively small constant $C>1$''. &
Low completeness, soundness smaller by the factor $q^{\beta}$. Not cited.\\
\addlinespace
O'Donnell and Wright, 2012 \cite{ODonnellWright2012-NewPointNPHardnessUniqueGames} &
``distinguishing $\tfrac12$-satisfiable Unique-Games instances from $(\tfrac38+\epsilon)$-satisfiable instances is NP-hard (for all $\epsilon>0$)''; ``the first hardness result showing that it helps to take the alphabet size larger than~2'' (abstract, p.~1); soundness $\tfrac38+q^{-\Theta(1)}$ for label size $q$ (Main Theorem, p.~3). &
$(1/2,\,3/8)$, alphabet growing with the soundness error. Not cited.\\
\addlinespace
H\aa stad, Huang, Manokaran, O'Donnell and Wright, 2015/2017 \cite{HastadEtAl2015-ImprovedNPInapprox2LinearEquations} &
``Fix any $C<11/8$. Then it is NP-hard to $(\epsilon,C\epsilon)$-approximate 2-Lin(2) (for any $0<\epsilon\le\epsilon_0=1/8$)'' (Thm~1.7, p.~4; Cor.~3.3, p.~18, also for Max-Cut); ``the best known NP-hardness even for the Unique-Games problem'' (abstract, p.~1). &
Binary alphabet, $1-\epsilon$ versus $1-C\epsilon$ for $C<11/8$. Not cited.\\
\addlinespace
Wiman, 2018 \cite{NL-Wiman2018} &
For $\epsilon<1/32$, given a Max-Cut instance $I$ it is NP-hard to distinguish $\operatorname{val}(I)\ge1-\epsilon$ from $\operatorname{val}(I)\le1-1.45685\,\epsilon$ (Thm~3.18, p.~36); for Max-2-Lin(2) with the constant $1.4568$ (Thm~1.5, p.~3). &
Binary alphabet, $1-\epsilon$ versus $1-1.45685\,\epsilon$. Not cited.\\
\addlinespace
Bhangale and Khot, 2019/2022 \cite{BhangaleKhot2019-UGHardnessToNPHardnessByLosingHalf} &
For every $\delta>0$ there is $L$ such that, given a Unique Games instance $G=(A,B,E,[L],\{\pi_e\})$ ``regular on the $A$ side, it is NP-hard to distinguish'' the YES case, a set $A'\subseteq A$ of size $(\tfrac12-\delta)|A|$ and an assignment satisfying all edges incident on $A'$, from the NO case, every assignment satisfies at most a $\delta$ fraction of the edges (Thm~1.2, p.~3). &
Completeness $1/2$ in the strong-left form of \ploc{\S8.1, p.~40}{07-consequences.tex:49--54} and \ploc{Lemma~A.1, p.~43}{finite-consequences.tex:22--34}, where the preprint has $1-\zeta$ for $\tfrac12-\delta$. Not cited.\\
\addlinespace
Martinsson, 2024 \cite{Martinsson2024-NPHardnessApproximationCurveMax2Lin2} &
For the curve $s(c):[0.5,1]\to[0.5,1]$ of Definition~40 and every small $\epsilon>0$, ``it is NP-hard to distinguish between instances of Max-2Lin(2)'' of value at least $c-\epsilon$ and of value at most $s(c)+\epsilon$ (Thm~1, p.~5); the curve ``either matches or improves all of the previously known inapproximability NP-hardness results for Max-2Lin(2)'' (abstract). &
Binary alphabet; a curve of points, with $(1-s(c))/(1-c)>1.48969$ for $c\ge0.9232$ (Cor.~2, p.~5). Not cited.\\
\addlinespace
Raghavendra and Steurer, 2010 \cite{RaghavendraSteurer2010-GraphExpansionUGC} &
Conjecture~1.3: for every $\eta>0$ there is $\delta$ such that distinguishing $\Phi_G(\delta)\ge1-\eta$ from $\Phi_G(\delta)\le\eta$ (Problem~1, PDF p.~2) is NP-hard; ``Theorem~1.4. The Gap-Small-Set Expansion conjecture implies the Unique Games Conjecture'' (PDF p.~3). Recorded in Khot's survey, which the preprint cites for the conjecture's formulation (\ploc{\S1.1, p.~2}{00-introduction.tex:41--42}), as \sloc{Khot2010-OnTheUGC}{Thm~5.2, p.~19}. &
The full conjecture, under the Small-Set Expansion hypothesis. Not cited; the preprint does not mention small-set expansion, its nearest passage being the sentence on Bafna et al.'s hypercontractivity certificates (\ploc{\S1.2, p.~3}{00-introduction.tex:103--108}).\\
\addlinespace
Eldan and Moshkovitz, 2020 \cite{EldanMoshkovitz2020-NonUniqueToBooleanUniqueGames} &
``Assume the Subspaces Near-Intersection Conjecture (Conjecture~1.7 in the sequel). For any $C\ge1$, for any sufficiently small $\delta>0$, distinguishing $1-\delta$ vs.\ $1-C\delta$ Boolean unique games is NP-hard'' (Thm~1.4, p.~2), the ``Boolean Unique Games Conjecture'' of Conj.~1.3, p.~2. &
Binary alphabet, $1-\delta$ versus $1-C\delta$ for every fixed $C$, under the Subspaces Near-Intersection hypothesis. Not cited.\\
\addlinespace
\bottomrule
\end{xltabular}
\end{center}
The one earlier point the preprint does name is credited:
\ploc{\S1.3, p.~3}{00-introduction.tex:129--131} attributes ``a Unique Games
gap with completeness close to one half and arbitrarily small soundness'' to
the 2-to-2 line it cites, which is Khot, Minzer and Safra's statement that
Gap UG$[\mathbb F_2^{\ell}](\tfrac12,\epsilon)$ is NP-hard for every
$\epsilon>0$ and large enough $\ell=\ell(\epsilon)$
\sloc{KhotMinzerSafra2018-PseudorandomGrassmannExpansion}{Thm~1.5, p.~6},
preceded by the $1/2$ versus $1/4+\epsilon$ point of Dinur, Khot, Kindler,
Minzer and Safra \sloc{DinurEtAl2018-Towards2to1}{Cor.~1.19, p.~15}.

\paragraph{How to check.}
\begin{enumerate}
\item \ploc{Theorem~1.1 (p.~2)}{00-introduction.tex:22--39} and the survey
  at \ploc{\S\S1.2--1.3, pp.~2--3}{00-introduction.tex:57--131} against
  \sloc{ODonnellWright2012-NewPointNPHardnessUniqueGames}{Main Theorem, p.~3}
  and \sloc{HastadEtAl2015-ImprovedNPInapprox2LinearEquations}{Thm~1.7, p.~4}:
  neither work, nor any other in the table, appears in the preprint's
  reference list, and the citations of H\aa stad's paper concern the
  3-variable parity gap and its Fourier decoding
  (\ploc{\S1.2, p.~2}{00-introduction.tex:72};
  \ploc{Theorem~4.1, p.~18}{03-game.tex:25--40};
  \ploc{\S5.6, p.~30}{04-matrix.tex:621--623}).
\item \ploc{\S8.1, p.~40}{07-consequences.tex:49--54} and
  \ploc{Lemma~A.1, p.~43}{finite-consequences.tex:22--34} against
  \sloc{BhangaleKhot2019-UGHardnessToNPHardnessByLosingHalf}{Thm~1.2, p.~3}:
  the same left-regular, strongly-satisfied-left-vertices form, at
  completeness $1-\zeta$ and $\tfrac12-\delta$ respectively.
\item \ploc{\S1.1, p.~2}{00-introduction.tex:41--42} and
  \ploc{\S1.2, p.~3}{00-introduction.tex:103--108} against
  \sloc{RaghavendraSteurer2010-GraphExpansionUGC}{Thm~1.4, PDF p.~3} and
  \sloc{Khot2010-OnTheUGC}{Thm~5.2, p.~19}: the preprint cites Khot's
  survey only for Conjecture~2.5 and does not mention small-set expansion.
\end{enumerate}

\space%
\finding[Earlier unconditional ordering bound]{Earlier unconditional ordering hardness from strong-left Unique Games}{strong-left-ordering-composition}

\paragraph{Summary.}
Bhangale and Khot \cite{BhangaleKhot2019-UGHardnessToNPHardnessByLosingHalf}
(CCC~2019; Theory of Computing 2022) prove that Maximum Acyclic Subgraph is
NP-hard to approximate within $2/3+\varepsilon$
(\sloc{BhangaleKhot2019-UGHardnessToNPHardnessByLosingHalf}{Thm~1.4, p.~5}),
by running the ordering reduction of Guruswami, Manokaran and Raghavendra
\cite{GuruswamiManokaranRaghavendra2008-BeatingRandomOrderingMAS} on a
left-regular Unique Games theorem in \emph{strong-left} form, where one
labeling satisfies every edge at a $1/2-\delta$ fraction of the left vertices
(\sloc{BhangaleKhot2019-UGHardnessToNPHardnessByLosingHalf}{Thm~1.2, p.~3};
\sloc{BhangaleKhot2019-UGHardnessToNPHardnessByLosingHalf}{\S5, pp.~21--23}).
The preprint does not cite Bhangale and Khot anywhere: the paragraph stating
the unconditional Maximum Acyclic Subgraph threshold $1/2$
(\ploc{\S8.2, p.~42}{07-consequences.tex:220--224}) does not record the
earlier unconditional bound, and
\ploc{Lemma~A.1, p.~43}{finite-consequences.tex:22--34} states Theorem~1.2's
strong-left form with $1-\zeta$ in place of $1/2-\delta$ (as a hardness point
for Unique Games, Theorem~1.2 is in finding~\ref{f:earlier-unique-games-hardness-points}).
The preprint derives its thresholds in
\ploc{\S8.2 (pp.~41--42)}{07-consequences.tex:160--230} and
\ploc{Appendix~A (pp.~43--54)}{finite-consequences.tex:1--988}: Lemma~A.1,
whose proof uses Khot and Regev's rounding device, credited there
(\ploc{p.~44}{finite-consequences.tex:36--39}), is fed into the reduction of
Guruswami, H\aa stad, Manokaran, Raghavendra and Charikar
\sloc{GuruswamiEtAl2011-BeatingRandomOrderingIsHard}{Conj.~3.8, p.~14; Thm~7.1, p.~20},
credited at \ploc{Proposition~A.2, p.~46}{finite-consequences.tex:242--265},
giving $1-\varepsilon$ versus $\rho_\Pi+\varepsilon$ for every ordering
predicate $\Pi$ (\ploc{Corollary~8.2, pp.~41--42}{07-consequences.tex:190--218}).
\rel{weaker result}: Theorem~1.4 is the earlier unconditional version of that
threshold, for Maximum Acyclic Subgraph only and at $2/3+\varepsilon$ instead
of $1/2+\varepsilon$, by the same composition, which both papers owe to
Guruswami et al.

\paragraph{Correspondence.}
\begin{center}\small
\begin{xltabular}{\linewidth}{@{}>{\raggedright\arraybackslash}p{0.15\linewidth}>{\raggedright\arraybackslash}X>{\raggedright\arraybackslash}X@{}}
\toprule
 & Preprint & Bhangale--Khot \cite{BhangaleKhot2019-UGHardnessToNPHardnessByLosingHalf}\\
\midrule
\endfirsthead
\toprule
 & Preprint & Bhangale--Khot \cite{BhangaleKhot2019-UGHardnessToNPHardnessByLosingHalf}\\
\midrule
\endhead
Input & \ploc{Lemma~A.1, p.~43}{finite-consequences.tex:22--34}: simple bipartite $\Phi_\varphi=(A,B,E)$ over $[R]$, left degree $h$; YES: one labeling strongly satisfies $\ge(1-\zeta)|A|$; NO: $\operatorname{val}<\nu$ & Thm~1.2 (p.~3): $G$ regular on the $A$ side; YES: $A'\subseteq A$, $|A'|=(1/2-\delta)|A|$, all incident edges satisfied; NO: $\operatorname{val}(G)\le\delta$\\
\addlinespace
Composition & $a$ uniform in $A$, neighbors $b_i$, positions $(b_i,\pi_{a\to b_i}\widetilde Z_i)$; order $f(b,z)=z_{L(b)}$, with $f(b,\pi_{a\to b}z)=z_{L(a)}$ at strongly satisfied $a$ (\ploc{(A.14)--(A.15), p.~51}{finite-consequences.tex:669--705}) & Figure~4 (p.~22): $u$ uniform in $A$, neighbors $v_1,v_2$, edge $(v_1,\tilde x\circ\pi_1)\to(v_2,\tilde y\circ\pi_2)$; order $O(v,x)=x_{\ell(v)}$, ``if $u\in A'$ then $\pi_1(\ell(v_1))=\pi_2(\ell(v_2))=\ell(u)$'' (Lemma~5.3, pp.~21--22)\\
\addlinespace
Result & \ploc{Corollary~8.2, pp.~41--42}{07-consequences.tex:190--218}: $1-\varepsilon$ vs.\ $\rho_\Pi+\varepsilon$ for every $\Pi$; Maximum Acyclic Subgraph threshold $1/2$ (\ploc{p.~42}{07-consequences.tex:220--224}) & Thm~1.4 (p.~5): Maximum Acyclic Subgraph NP-hard within $2/3+\varepsilon$; proof (p.~23) with $c=1/2-\delta$, completeness about $3/4$ against soundness $1/2+\varepsilon'$\\
\addlinespace
\bottomrule
\end{xltabular}
\end{center}

\paragraph{How to check.}
\begin{enumerate}
\item \sloc{BhangaleKhot2019-UGHardnessToNPHardnessByLosingHalf}{Thm~1.4, p.~5, and its proof, p.~23}
against \ploc{Corollary~8.2, pp.~41--42}{07-consequences.tex:190--218} and
\ploc{p.~42}{07-consequences.tex:220--224}: $2/3+\varepsilon$ and
$1/2+\varepsilon$ for the same problem; the earlier bound is not mentioned.
\item \sloc{BhangaleKhot2019-UGHardnessToNPHardnessByLosingHalf}{Thm~1.2, p.~3}
against \ploc{\S8.1, p.~40}{07-consequences.tex:49--54} and
\ploc{Lemma~A.1, p.~43}{finite-consequences.tex:22--34}: the same
left-regular strong-left form at completeness $1/2-\delta$ and $1-\zeta$.
\item \sloc{BhangaleKhot2019-UGHardnessToNPHardnessByLosingHalf}{Figure~4, p.~22, and Lemma~5.3, pp.~21--22}
against \ploc{(A.14)--(A.15), p.~51}{finite-consequences.tex:669--705}, both
against
\sloc{GuruswamiEtAl2011-BeatingRandomOrderingIsHard}{Conj.~3.8, p.~14; \S7.1, pp.~20--21}:
both run the composition proposed there.
\end{enumerate}

\space%

\begin{thebibliography}{99}
\small
\bibitem{DinurEtAl2016-Towards2to1-ECCC} Irit Dinur, Subhash Khot, Guy Kindler, Dor Minzer, Muli Safra, Towards a Proof of the 2-to-1 Games Conjecture?, ECCC TR16-198, December 14, 2016. ECCC TR16-198.
\bibitem{DinurEtAl2018-Towards2to1} Irit Dinur, Subhash Khot, Guy Kindler, Dor Minzer, Muli Safra, Towards a Proof of the 2-to-1 Games Conjecture?, Theory of Computing 21 (2025) 1--50; STOC 2018. doi:10.4086/toc.2025.v021a011.
\bibitem{KhotMinzerSafra2017-IndependentSets2to2Grassmann} Subhash Khot, Dor Minzer, Muli Safra, On Independent Sets, 2-to-2 Games and Grassmann Graphs, Theory of Computing 21 (2025) 1--55; STOC 2017. doi:10.4086/toc.2025.v021a010.
\bibitem{KhotSafra2013-TwoProverStrongSoundness} Subhash Khot, Muli Safra, A Two-Prover One-Round Game with Strong Soundness, Theory of Computing 9 (2013) 863--887; FOCS 2011. doi:10.4086/toc.2013.v009a028.
\bibitem{BellareGoldreichSudan1998-FreeBitsPCPs} Mihir Bellare, Oded Goldreich, Madhu Sudan, Free Bits, PCPs and Non-Approximability -- Towards Tight Results, SIAM Journal on Computing 27 (1998) 804--915. doi:10.1137/S0097539796302531.
\bibitem{KhotNaor2008-ApproximateKernelClustering} Subhash Khot, Assaf Naor, Approximate kernel clustering, Mathematika 55 (2009) 129--165; FOCS 2008; arXiv v1 July 29, 2008, v2 December 9, 2008. arXiv:0807.4626v2; doi:10.1112/S002557930000098X.
\bibitem{ChawlaEtAl2006-HardnessMulticutSparsestCut} Shuchi Chawla, Robert Krauthgamer, Ravi Kumar, Yuval Rabani, D. Sivakumar, On the Hardness of Approximating Multicut and Sparsest-Cut, Computational Complexity 15(2) (2006) 94--114; CCC 2005. doi:10.1007/s00037-006-0210-9.
\bibitem{DemaineEtAl2006-CorrelationClusteringGeneralWeightedGraphs} Erik D. Demaine, Dotan Emanuel, Amos Fiat, Nicole Immorlica, Correlation Clustering in General Weighted Graphs, Theoretical Computer Science 361(2--3) (2006) 172--187. doi:10.1016/j.tcs.2006.05.008.
\bibitem{ODonnellWright2012-NewPointNPHardnessUniqueGames} Ryan O'Donnell, John Wright, A new point of NP-hardness for Unique Games, STOC 2012; author manuscript dated September 30, 2012.
\bibitem{HastadEtAl2015-ImprovedNPInapprox2LinearEquations} Johan Håstad, Sangxia Huang, Rajsekar Manokaran, Ryan O'Donnell, John Wright, Improved NP-Inapproximability for 2-Variable Linear Equations, Theory of Computing 13 (2017) 1--51; APPROX 2015. doi:10.4086/toc.2017.v013a019.
\bibitem{BhangaleKhot2019-UGHardnessToNPHardnessByLosingHalf} Amey Bhangale, Subhash Khot, UG-hardness to NP-hardness by Losing Half, Theory of Computing 18(5) (2022) 1--28; CCC 2019. doi:10.4086/toc.2022.v018a005.
\bibitem{Martinsson2024-NPHardnessApproximationCurveMax2Lin2} Björn Martinsson, On the NP-Hardness Approximation Curve for Max-2Lin(2), arXiv v1 August 9, 2024, v2 August 12, 2024 (manuscript dated August 13, 2024). arXiv:2408.04832v2.
\bibitem{RaghavendraSteurer2010-GraphExpansionUGC} Prasad Raghavendra, David Steurer, Graph Expansion and the Unique Games Conjecture, STOC 2010.
\bibitem{EldanMoshkovitz2020-NonUniqueToBooleanUniqueGames} Ronen Eldan, Dana Moshkovitz, Reduction From Non-Unique Games To Boolean Unique Games, arXiv v1 June 23, 2020, v2 July 8, 2021. arXiv:2006.13073v2.
\bibitem{NL-Wiman2018} M{\aa}rten Wiman, Improved Inapproximability of Max-Cut through Min-Cut, Master's thesis, KTH Royal Institute of Technology, School of Electrical Engineering and Computer Science, Stockholm (June 15, 2018), DiVA record diva2:1229467
\bibitem{NL-FeigeReichman2004} Uriel Feige, Daniel Reichman, On Systems of Linear Equations with Two Variables per Equation, APPROX-RANDOM 2004, Lecture Notes in Computer Science 3122, Springer (2004), 117--127, doi:10.1007/978-3-540-27821-4\_11
\bibitem{Hastad2001-SomeOptimalInapproximability} Johan Håstad, Some optimal inapproximability results, Journal of the ACM 48 (2001) 798--859; author manuscript dated February 18, 2002. doi:10.1145/502090.502098.
\bibitem{OAI-Preprint} OpenAI, \emph{The Unique Games Theorem}, OpenAI Math Release preprint, September 23, 2026, \url{https://github.com/openai/math}.
\bibitem{OAI-Companion1} OpenAI, \emph{Perfect completeness for 2-to-1 games}, OpenAI Math Release preprint, September 23, 2026 (cited by the preprint as \texttt{OpenAIPerfectTwoToOne2026}).
\bibitem{OAI-Companion2} OpenAI, \emph{Constant-factor hardness of uniform sparsest cut}, OpenAI Math Release preprint, September 24, 2026 (cited by the preprint as \texttt{OpenAIUniformSparsest2026}).
\bibitem{OAI-Companion3} OpenAI, \emph{The Gaussian propeller bound in every dimension}, OpenAI Math Release preprint, September 24, 2026 (cited by the preprint as \texttt{OpenAIGaussianPropeller2026}).
\bibitem{KhotEtAl2005-OptimalInapproxMaxCut} Subhash Khot, Guy Kindler, Elchanan Mossel, Ryan O'Donnell, Optimal Inapproximability Results for MAX-CUT and Other 2-Variable CSPs?, FOCS 2004; SIAM Journal on Computing 37 (2007) 319--357; author manuscript ECCC TR05-101 (September 2005). doi:10.1137/S0097539705447372; ECCC TR05-101.
\bibitem{KhotMinzerSafra2018-PseudorandomGrassmannExpansion} Subhash Khot, Dor Minzer, Muli Safra, Pseudorandom Sets in Grassmann Graph have Near-Perfect Expansion, ECCC TR18-006 (January 2018), revision 2 dated May 19, 2018; FOCS 2018; Annals of Mathematics 198 (2023) 1--92. ECCC TR18-006 rev. 2; doi:10.4007/annals.2023.198.1.1.
\bibitem{Khot2010-OnTheUGC} Subhash Khot, On the Unique Games Conjecture, CCC 2010 (invited survey), 99--121. doi:10.1109/CCC.2010.19.
\bibitem{GuruswamiManokaranRaghavendra2008-BeatingRandomOrderingMAS} Venkatesan Guruswami, Rajsekar Manokaran, Prasad Raghavendra, Beating the Random Ordering is Hard: Inapproximability of Maximum Acyclic Subgraph, FOCS 2008, 573--582.
\bibitem{GuruswamiEtAl2011-BeatingRandomOrderingIsHard} Venkatesan Guruswami, Johan Håstad, Rajsekar Manokaran, Prasad Raghavendra, Moses Charikar, Beating the Random Ordering is Hard: Every ordering CSP is approximation resistant, SIAM Journal on Computing 40(3) (2011) 878--914; ECCC TR11-027 (2011); preliminary versions FOCS 2008 and CCC 2009. ECCC TR11-027; doi:10.1137/090756144.
\end{thebibliography}
\space%

\end{document}
